% 本文件由 LaTeX 排版 Agent 根据有效 translation.json 自动生成。
% author=John Chrysostom; work_id=1902; title=不伤害自己的人，无人能伤害他

\chapter{不伤害自己的人，无人能伤害他}

% structure: p0001 source=h1 target=omitted reason=page-title-already-rendered-by-parent

我深知，对于那些心性粗鄙、贪恋现世事物、拘泥尘世、受肉体享乐奴役、对精神理念毫无坚定把握之人，这篇论述将是陌生而悖理的。他们必会纵声嘲笑，并从一开始就指责我说出难以置信之事。然而，我不会因此而放弃我的承诺；正因如此，我将以极大的热忱着手证明我所承担之事。因为，若持此看法之人愿意不再喧哗骚动，而是等候我言论的终结，我确信他们将站到我一边，自责之前所受的欺骗，并收回己见、表示歉意，恳请宽恕他们对此事所持的错误观点，并对我深表感激——正如病患在摆脱缠身的身体疾苦后，对医师所怀的感激那样。请勿向我提及你眼下心中流行的判断，且待聆听我论证的争辩，然后你便能作出公正的裁决，而不因无知而阻碍你形成真确的判断。因为即使在世俗案件中，法官若见第一位演说家口若悬河、滔滔不绝，以言辞压倒一切，他们也不敢在尚未耐心听取对立辩手的反驳前作出判决；即使第一位演说家的言论看似完全公正，他们也为第二位保留无偏见的听证。事实上，法官的独特功绩在于尽可能准确地查明每一方的主张，然后提出自己的判断。\par

现在，代替演说家的，是人类的普遍假设——它历经世代，在众人心中深植，并在此世间如此宣称：它说\Quote{万物}都已颠倒，人类充满混乱，每日有许多人受委屈、受侮辱、遭暴力与伤害——弱者受强者欺，穷人受富人凌；而且，正如无法数清海中的波浪，同样无法计谋那些受阴谋、侮辱与苦难的人数；法律的矫正、审判的恐惧，或任何其他事物，都无法遏止这场瘟疫与混乱，反而罪恶日增，受难者的呻吟、哀叹与哭泣遍及各处。而受命纠正此类恶行的法官，他们自己反而加剧了风暴，点燃了混乱。因此，许多更为愚钝可鄙之人，被一种新的疯狂攫住，竟责怪天主的眷顾——当他们看见忍耐之人常被强夺、折磨、压迫，而狂妄、急躁、低贱、出身卑微之人却致富、掌权、令众人畏惧，并给温和之人带来无数烦扰，且此事行于城镇、乡村、旷野、海陆各处。我们这篇论述必然与所声称的针锋相对，持有一个新的——如我起初所言——与意见相左，却有益而真实，并有益于那些留意并信服之人。因为我所承担的是证明（只是请不要喧哗）：没有一个受委屈的人是受他人委屈，而是自己经历这种伤害。\par

2. 但为使我的论证更明晰，让我们首先探讨不义是什么，它通常由哪些事物构成；以及人的德行是什么，什么能毁坏它；此外，什么看起来能毁坏它而实际上不能。例如（因为我必须借助例子来完成我的论证），每种事物都受一种能毁坏它的恶所支配：铁受锈，羊毛受蛀虫，羊群受狼。葡萄酒的德行在发酵变酸时受损；蜂蜜在失去天然甜味而变成苦液时受损。谷穗因霉病和干旱而毁坏，葡萄的果实、叶子和枝条因有害的蝗虫群而毁坏，其他树木因毛虫而毁坏，无理性的受造物因各种疾病而毁坏；为了不因遍历所有可能的例子而延长清单，我们自己的肉身也易受热病、瘫痪和许多其他疾病的侵害。既然这些东西每一件都容易受其德行毁坏之物的影响，现在让我们考虑什么伤害人类，什么毁坏人的德行。大多数人认为有许多不同的事物有这种效果；因为我必须提及关于此事的错误意见，并在驳斥它们之后，进而展示真正毁坏我们德行的是什么：并清楚证明，除非我们背叛自己，否则无人能施加这种伤害或带来这种毁坏。于是，持有错误意见的群众想象有许多不同的事物毁坏我们的德行：有些人说是贫穷，有些人说是身体疾病，有些人说是财产损失，有些人说是诽谤，有些人说是死亡；他们不断地哀叹痛哭这些事：而当他们怜悯受苦者、流泪时，他们激动地相互喊道：\Quote{某某人遭遇了何等灾难！他一瞬间被夺走了全部财产。}另一个人又会说：\Quote{某某人患了重病，被在场的医生放弃了希望。}有些人哀叹痛哭监狱里的囚犯，有些人哀叹那些被逐出本国、流放到流放地的人，其他人哀叹那些被剥夺自由的人，其他人哀叹那些被敌人抓获成为俘虏的人，其他人哀叹那些被淹死、烧死或被房屋倒塌掩埋的人，但没有人哀叹生活在邪恶中的人：相反，比一切都更糟的是，他们常常祝贺这些人，而这是万恶之源。来吧（只希望你按我一开始的劝告，不要骚动），让我证明，上述事物没有一个伤害那节制生活的人，也不能毁坏他的德行。因为请告诉我，如果一个人失去了一切，或因诽谤者，或因强盗，或因狡猾的仆人被剥夺了财产，这种损失对人的德行有何伤害？\par

但如果可行，让我先指出人的德行是什么，从处理另一类存在物的事例开始，以便使大多数读者更容易理解和清楚。\par

3. 那么，马的美德是什么？是有一条镶金之缰绳，与相配的肚带，以及以丝线制成的、用来固定鞍垫的带子，并有各色与金线织成的衣服，饰有珠宝的头饰，以及由金线编成的鬃毛？还是它奔跑迅速、腿力强劲、步伐稳健、蹄子适合良驹，并有勇气胜任长途跋涉与征战，在战场上能保持冷静，若发生溃逃则能拯救其骑手？难道不明显这些才是构成马的美德的东西，而非其它？再说，你认为驴和骡子的美德是什么？岂不是忍耐负重、轻松完成旅程、蹄子坚如岩石？我们能说它们的外在装束对其自身的美德有任何贡献吗？绝不。我们又会称赞哪种葡萄树？是枝叶繁茂的，还是果实累累的？我们又给橄榄树下什么美德的判断？是拥有大树枝和极旺盛的叶子，还是展示其专有之果实遍满全树？那么，让我们对人也如此行：让我们确定人的美德是什么，并且只将那种毁灭美德之事视为伤害。那么，人的美德是什么？不是财富，以致你该畏惧贫穷；也不是身体健康，以致你该害怕疾病；也不是公众的意见，以致你该对恶名感到惊慌；也不是纯粹为生命本身而活，以致死亡对你可怕；也不是自由，以致你该逃避奴役；而是持守真道之谨慎，与生活之正直。这些事，连魔鬼本身也不能从人那里夺走，如果拥有它们的人以必要之谨慎守护它们：那个最恶毒凶残的恶魔知道此事。为此，他也夺去了约伯的财产，不是要使他贫穷，而是要逼他说出亵渎之言；他也折磨他的身体，不是要使他患病，而是动摇他灵魂的美德。但尽管如此，当他发动所有计谋，使他从富人变成穷人（那在我们看来最可怕的灾祸），使那曾儿女成群的人变成无子嗣，使他的全身比公堂上的刽子手更残酷地折磨（因为他们的指甲撕裂那些落入他们手中之人的肋骨，不如虫蛀撕裂他的身体那么厉害），并在他身上附上恶名（因为约伯在场的朋友们说：\Quote{你并没有受到与你罪恶相称的惩罚，}并说了许多指责他的话），并且不仅将他逐出城市与家庭，转移到另一个城市，而且实际上使粪堆成为他的家与城市；在此之后，他不仅没有对他造成损害，反而因他所设的计谋使他更加荣耀。他不仅没有夺走他的任何财产，尽管夺走了这么多，反而增加了他美德的财富。因为在此之后，他享有更大的信心，因为他进行了更严峻的竞赛。如果那忍受如此苦难之人，而且并非由人手中，而是由比所有人更邪恶的魔鬼手中所施，没有遭受伤害，那么那些说某某人伤害了我、损坏了我的人，将来还有什么辩护？因为如果那充满如此巨大恶意的魔鬼，在发动所有工具、放出所有武器、将一切加诸于人的灾祸以极致程度倾泻在那义人的家庭与本人身上之后，却并没有伤害他，反而如我所说使他获益：那么某些人怎能指责某某人，声称自己遭受了他们的伤害，而非自己的呢？\par

4. 那么，怎么办？有人会说：他不是伤害了亚当，使他跌倒，并把他逐出乐园吗？不：不是他做的，而是那受害者的懈怠，以及他缺乏节制与警醒，才是原因。因为那运用如此强大而多样的手段，却无法制伏约伯的人，又怎能用次等的手段制伏亚当，除非亚当由于自己的懈怠而背叛了自己？那么，那被诽谤、财物被没收、失去所有财产、被逐出祖业、挣扎于极度贫困中的人，难道没有受到伤害吗？不！他没有受害，甚至受益了，只要他保持清醒。因为，请告诉我，这给宗徒们带来了什么害处？他们不是一直在忍受饥饿、干渴和赤身露体吗？而这正是他们如此卓越、杰出，并从天主那里获得许多帮助的原因。再问：拉匝禄的疾病、疮口、贫困和无人保护，对他造成了什么损害？难道不是这些为他编织了更丰盛的胜利花冠吗？或者，若瑟在自己乡里和异乡都受恶名——他被当成奸夫和淫乱者——这给他带来了什么害处？为奴或离乡背井，又对他有何损害？不正是因了这些事，我们才对他充满钦佩和惊叹吗？我何必谈论流放异乡、贫困、恶名和奴役呢？因为死亡本身对亚伯尔造成了什么伤害？尽管那是暴虐的、过早的死亡，而且是出于兄弟之手？这不正是他的赞美传遍天下的原因吗？你看，这篇讲论所证明的，比它所应许的还多！因为它不仅揭示了没有人会受任何人的伤害，而且还指出那些留心自己的人，从（这样的攻击）中获得更大的益处。那么，有人会说，惩罚和刑罚的目的是什么？地狱的目的是什么？如果没有人受伤害，也没有人伤害别人，那么这许多威胁又有什么意义？你说什么？你为何混淆论点？因为我没有说没有人伤害，而是说没有人受伤害。你怎么可能，你会说，当许多人进行伤害时，却没有人受伤害？正如我刚才所指出的那样。若瑟的兄弟们确实伤害了他，但他自己并没有受害；加因算计了亚伯尔，但他自己并没有被算计。这就是为何有惩罚和刑罚。因为天主并没有因受苦者的德行而废除惩罚；而是因作恶者的恶意而规定刑罚。因为虽然受虐待的人因针对他们的阴谋而变得更加卓越，但这并不是策划阴谋者的意图所致，而是受害者勇气的结果。因此，对后者而言，智慧的赏报已经预备好；对前者而言，邪恶的惩罚也已备妥。你被剥夺了钱财吗？请读这句经文：\BibleQuote{我赤身出离母胎，也要赤身归去。}\BibleRef{约 1:21}。再加上这句宗徒的话：\BibleQuote{因为我们没有带什么到世界上来，也不能带什么去。}\BibleRef{弟前 6:7}。你被恶名中伤，有人用无数辱骂攻击你吗？记住那段经文：\BibleQuote{几时众人都夸赞你们，你们是有祸的！}\BibleRef{路 6:26}，以及\Quote{当他们因你们的名而辱骂你们时，你们欢喜踊跃吧！}。你被流放到异地了吗？要想到你在此没有祖国，而且你若明智，就应把整个世界视为异乡。或者你染上了重病吗？请引用宗徒的话：\BibleQuote{我们外在的人日渐朽坏，内在的人却日日更新。}\BibleRef{格后 4:16}。有人死于非命吗？请考虑若翰的情况：他在监牢里被斩首，放在盘子里，成了舞女的赏报。思考从这些事中所产生的回报：因为所有这些痛苦，当被任何人加于他人身上而不义时，便赎罪并成就义德。对于那些勇敢承受它们的人而言，这些痛苦的好处是如此之大。\par

5. 既然金钱的损失，毁谤，辱骂，放逐，疾病，酷刑，以及那似乎比这一切更可怕的事，即死亡，都不能伤害那些忍受它们的人，反而增加了他们的益处，那么你从何能向我证明，一个人在这些事上根本没有受到伤害时，他却是受了伤害呢？因为我将努力证明相反的事，指出那些受到最大伤害和侮辱，并忍受最不可治愈的恶事的人，正是做这些事的人。因为有什么比\Person{加音}的状况更悲惨的呢？他这样对待他的兄弟。有什么比\Person{斐理伯}的妻子更可怜的呢？她斩了\Person{若翰}的头。或者是\Person{若瑟}的兄弟们，他们把他卖掉，把他送到流徙之地。或者是魔鬼，它用如此巨大的灾难折磨\Person{约伯}。因为它不仅为了其他的不义，同时也为了这次攻击，将付出不小的代价。你看到吗？这里的论证甚至比所提出的还要多，表明那些受侮辱的人不仅在这些攻击中没有受到伤害，而且所有的祸害都回旋到策划它们的人的头上。因为既然人的德性不在于财富、自由、故乡的生活，也不在于我所提到的其他东西，而只在于灵魂的正直行为，那么当伤害针对这些事时，人的德性本身自然不会受到伤害。那么怎么样？假如有人伤害了灵魂的道德状态呢？即使在那时，如果一个人受到损害，损害也不是来自他人，而是来自内部，来自他自己。\Quote{何以见得？}你说。当一个人被另一个人殴打，或者被剥夺了他的财产，或者忍受了其他严重的侮辱，说出了亵渎的话，他当然因此受到了损害，而且是非常大的损害，然而它并非来自施加侮辱的人，而是来自他自己灵魂的狭隘。因为我将重复我以前说过的话，没有人即使他极其邪恶，能够比那个对我们怀有刻骨仇恨的报复性的魔鬼更恶毒、更猛烈地攻击人：但这个残忍的魔鬼却没有能力动摇或推翻那个在法律之前、在恩宠时代之前生活的人，尽管它从四面八方向他发射了如此多而猛烈的武器。这就是灵魂高贵的力量。我还能说什么关于\Person{保禄}的呢？他难道没有忍受如此多的痛苦，以至于连列举它们都不容易？他被下狱，戴锁链，被拖来拖去，被犹太人鞭打，被石头砸，背上不仅被皮条刮伤，还被棍棒打伤，他被投入海中，屡次被强盗围困，与同胞发生冲突，不断受到敌人和熟人的攻击，遭受无数阴谋，与饥饿和赤身搏斗，经历其他频繁而持久的灾祸和患难：我还需要提到它们的大部分吗？他天天死：但是，尽管遭受了如此多而沉重的苦难，他不仅没有说出亵渎的话，反而因这些事喜乐，并以此夸耀：有一次他说：\BibleQuote{我在苦难中喜乐，}\BibleRef{哥 1:24} 接着又说：\BibleQuote{不但如此，我们在磨难中也夸耀。}\BibleRef{罗 5:3} 如果他因遭受如此大的困难而喜乐并夸耀，那么如果你在连它们极小一部分都没有遭受的情况下就亵渎，你将有什么借口，又有什么辩护呢？\par

6. 但有人会说：\Quote{我是在其他方面受到损害，即使我不亵渎，然而当我的钱财被抢走时，我就无法施舍了。}这不过是借口和托词。因为如果你为此而悲伤，你要确知贫穷并不妨碍施舍，即使你极端贫穷，你也不会比那位仅有一把面粉的妇人\BibleRef{列上 17:12}和那位只有两个小钱的妇人\BibleRef{路 21:2}更贫穷；她们各人都将全部所有给了那些需要的人，因而备受赞叹：如此巨大的贫穷并未阻碍如此伟大的仁爱，相反，从这两个小钱所出的施舍是如此丰盛和慷慨，足以使所有富人黯然失色，并且以心意之富和热忱之丰超越了那些投入大量钱币的人。因此，即使在这件事上你也没有受损，反而受益，借着微薄的奉献获得了比那些投入巨款的人更光荣的赏报。但是，即便我永远说这些话，那些沉溺于世俗事物、陶醉于眼前之乐的感官性人物也不会轻易放弃那些凋谢的花朵（因为今生的享乐正是如此）或放下其幻影：但更好的一类人的确同时紧握这两者，而更可怜和卑贱的一类人则更执着于前者而非后者。来吧，让我们剥去那些遮掩这些事物丑陋面貌的令人愉快的面具，揭露那娼妓的可憎面目。因为这种专注于奢华、财富和权力的生活正是如此：它是可憎、丑陋、充满可憎之事、令人不快、沉重且充满苦涩。这确实是这种生活的特点，它使被其迷惑之人失去一切借口：即，尽管它是他们渴望和努力的目标，却充满了许多烦恼和苦涩，蕴藏着无数邪恶、危险、流血、悬崖、峭壁、谋杀、恐惧与战栗、嫉妒与敌意、阴谋、不断的焦虑与挂虑，并且从这些巨大的邪恶中得不到任何益处，也不结出任何果实，除了惩罚和报复以及无尽的折磨。但尽管它具有这样的性质，对大多数人来说，它似乎仍是值得追求和激烈争夺的目标，这正表明被它迷惑之人的愚昧，而非事物本身的福乐。的确，小孩子对玩具热衷和兴奋，无法注意那些适合成年人事物。他们的不成熟情有可原：但这些人被剥夺了辩护的权利，因为他们虽已成年，却心性幼稚，在生活方式上比孩童更愚蠢。\par

现在告诉我，为何财富成为追求的目标？因为必须从这一点开始，因为对于大多数被这种严重疾病折磨的人来说，财富似乎比健康、生命、公众声誉、良好评价、国家、家庭、朋友和亲属以及其他一切更为宝贵。而且，火焰已升至云端：这炽热占据了陆地与海洋。也没有任何人能扑灭这火：所有人都忙着煽风点火，无论是已被火抓住的人，还是尚未被抓住的人，为了自己也被捕捉。你可以看到每个人——丈夫、妻子、家奴与自由人、富人与穷人——各按其能力，日夜背负着给这火添柴的负荷：不是木柴或柴捆（因为不是那种火），而是灵魂与身体的负荷，以及不义与邪恶。因为这种火正是以这样的材料点燃的。那些拥有财富的人，即使征服了整个世界，也从不给这可怕的激情设立任何界限：穷人也努力赶超他们，一种无法治愈的狂热、不可遏制的疯狂和不可救药的疾病占据了所有人的灵魂。而且这种情感战胜了其他一切，将其赶出灵魂：朋友和亲属都不被顾及；为何我提朋友和亲属？甚至妻子和儿女也不被考虑，还有什么比这些人更亲爱呢？但一切都摔在地上被践踏，当这野蛮而无人情味的女主人抓住了所有被她俘虏之人的灵魂时。因为她像一个无人情的女主人、严厉的暴君、野蛮的蛮族、公开而昂贵的娼妓，她用无数的危险和折磨贬低、耗尽、惩罚那些选择服从她的人；然而，尽管她可怕、严厉、凶残、残酷，有着蛮族的面孔，或者更确切地说，是野兽的面孔，比狼或狮子更凶猛，但被俘之人却认为她温柔可爱、比蜜更甜。尽管她每天为他们锻造刀剑和武器，挖陷阱，引他们到悬崖峭壁，编织无尽的惩罚罗网，她却使这些东西对被俘之人成为追求的目标，并让那些渴望被俘之人心向往之。正如猪喜欢并陶醉于在沟渠和泥沼中打滚，甲虫喜欢在粪便上不断爬行；同样，被贪财之欲俘虏的人比这些生物更可怜。因为这里的可憎之物更大，泥沼更令人作呕：他们沉溺于这种激情，想象从中获得许多快乐：这快乐并非来自事物本身的性质，而是来自被这种非理性品味折磨的理解力。这种品味在他们身上比在畜生身上更糟：因为正如泥沼和粪便，其快乐原因不在它们自身，而在那些陷入其中的生物的非理性本性；同样，你要认为人类也是如此。\par

7. 我们如何才能医治那些如此倾向的人？如果他们愿意听我们的话，敞开心扉，接受我们的言语，这是可能的。因为要将无理性的动物从它们不洁的习惯中扭转和转移是不可能的；它们缺乏理性。但是，在所有族类中最温和的、以理性和言语为荣的族类——我指的是人性——如果它愿意，就可以轻易而轻松地从泥沼、恶臭、粪堆及其可憎之物中解脱出来。因为，人啊，为什么财富在你眼中值得如此勤奋的追求？是因为桌上带来的快乐吗？还是因为那些奉承你的人因你的财富而给予的荣誉和随从？是因为你能够抵御那些烦扰你的人，并使所有人畏惧你？因为你不能提出任何其他原因，除了快乐、奉承、恐惧和报复的力量；因为财富通常不会使任何人更明智、更自制、更温和、更聪明、更仁慈、更慷慨、更不易发怒、更不贪食或更不贪乐；它不会训练任何人节制，或教导他如何谦卑，也不会在灵魂中引入或植入任何其他德行。你也不能说出这些事中的哪一件值得如此勤奋地寻求和渴望。因为财富不仅不知道如何种植和培育任何好东西，而且即使它找到了储存的好东西，它也会破坏、阻碍和摧残它们；甚至将它们连根拔起，并引入它们的对立面：无度的放纵、不合时宜的愤怒、不义的怒气、骄傲、傲慢、愚蠢。但让我不要说这些；因为那些被这疾病抓住的人不会忍受听到关于德行和恶习的事，他们完全沉溺于享乐，因此受制于它。那么现在让我们暂时搁置这些考虑，提出其余的问题，看看财富是否有任何快乐或荣誉：因为在我看来，情况恰恰相反。首先，如果你愿意，让我们调查富人和穷人的饮食，并问宾客们，哪一方享受最纯粹、最真实的快乐；是他们——整天斜靠在躺椅上，将早餐和正餐合并，使胃扩张，麻木感官，用过多的食物压沉船只，使船浸水，如同身体遭遇船难，制造镣铐、手铐和口塞，用醉酒和饱食的束缚——比铁链更沉重——捆绑整个身体，享受不到任何不受可怕梦境干扰的健全、纯净的睡眠，比疯子更悲惨，将某种自我施加的恶魔引入灵魂，在仆人眼中显现为笑柄，或者对其中较善良的人而言，成为一个引人流泪的悲剧景象，无法认出在场的人，不能说话或听见，必须被从躺椅上抬到床上——还是他们——清醒警觉，以需要为限进食，顺风航行，将饥饿和口渴视为食物和饮料的最佳调味品？因为没有什么比在进食时感到饥饿和口渴更能促进享受和健康，并以单纯的食物需要为满足，从不超越这个限度，也不将超出身体承受能力的负担强加于身体。\par

8. 如果你不相信我的话，就研究一下各阶层人的身体状况和灵魂。那些如此有节制地生活的人（因为不要告诉我那些罕见的情况，尽管有些人可能因其他原因而虚弱，但你要从那些常见的事例中判断），我说，他们的身体不是强健的吗？他们的感官不是清晰的吗？能够非常轻松地履行其适当的功能？而另一类人的身体则是松弛的，比蜡还柔软，并伴有多种疾病？因为痛风很快缠上他们，还有不合时宜的瘫痪、早衰、头痛、胀气、消化不良、食欲不振；他们需要不断看医生，长期服药，每日照料。这些事令人快乐吗？告诉我。那些真正知道快乐是什么的人会这样说吗？因为快乐是在欲望引导下产生的，随后得到满足：如果得到了满足，但欲望却不存在，快乐的条件就消失了。因此，病人即使面前摆着最迷人的食物，也会带着厌恶和压抑感尝一点：因为没有欲望给享受增添美味。因为产生欲望、并能够引起快乐的，不是食物或饮品的性质，而是进食者的胃口。因此，一位明智的人，对有关快乐的一切有准确了解，并知道如何关于这些事进行道德说教，曾说：\BibleQuote{饱足的灵魂嘲笑蜂房：} 表明快乐的条件不在于食物的性质，而在于进食者的心情。因此，先知在叙述埃及和旷野中的奇迹时，也提到了这一点：\BibleQuote{他以岩石中的蜂蜜满足了他们。} 然而，在任何地方都没有显示实际上有蜂蜜从岩石中为他们涌出：那么这句话是什么意思？因为百姓因劳累和长途跋涉而疲惫不堪，又因极度口渴而冲向清凉的泉水，他们对饮水的渴望作为调味品，作者为了描述他们从那些泉水中获得的快乐，称水为蜂蜜，并不是说元素变成了蜂蜜，而是说从水中获得的快乐可与蜂蜜的甜味相媲美，因为那些分享它的人急切地冲过去喝。\par

既然这些事如此，无论多么愚蠢的人也不能否认：难道不是非常明显，纯粹、不掺杂、活泼的快乐在穷人的餐桌上才能找到？而在富人的餐桌上，却有不适、厌恶和玷污？正如那位智者所说：\Quote{即使是甜美的事物也似乎是烦恼。}\par

9. 但是，有人会说，财富给拥有者带来荣誉，并使他们能轻易地向敌人报仇。请问，难道这就是财富之所以显得值得追求和争夺的原因——因为它们滋养了我们本性中最危险的情欲，助长愤怒付诸行动，膨胀野心空虚的泡沫，并刺激和驱使人们走向傲慢？这恰恰是我们理应坚决弃绝财富的理由，因为它们在我们心中引入了某些凶猛危险的野兽，剥夺了我们可能从所有人那里得到的真正荣誉，却向受骗者引入另一种与此相反的荣誉，只是涂上了它的色彩，并说服他们幻想这荣誉是相同的，而按其本性并非如此，只是表面上看起来如此。因为，正如娼妓的美色由染料和颜料装扮而成，缺乏真正的美丽，却使一张丑陋难看的脸在被它迷惑的人眼中显得美丽漂亮，而实际上并非如此：同样，财富也迫使谄媚看起来像是荣誉。因为我恳求你们不要考虑那些出于恐惧和奉承而公开表达的赞美：因为这些只是装饰和颜料；但请揭开每个如此奉承你的人的良心，在里面你将看到无数控诉者宣告反对你，比你的死敌和仇人更憎恶和厌恶你。如果有一天环境发生变化，揭开并暴露这张由恐惧制造的面具，正如太阳发出比平常更热的光线时，会揭露我提到过的那些女人的真实面容，那么你会清楚地看到，在整个过去的时间里，你被那些奉承你的人极其轻视，你以为自己从那些彻底恨你的人那里得到了荣誉，而他们心中却对你倾泻无尽的辱骂，渴望看到你陷入最大的灾难。因为没有什么比德行更能产生荣誉——这种荣誉既非强迫，也非伪装，亦非隐藏在欺骗的面具之下，而是真实而纯洁，能经受艰难岁月的考验。\par

10. 但你若想报复那些惹恼你的人吗？正如我刚才所说的，这正是财富特别应当被避免的原因。因为它促使你挥剑自戕，使你在未来的审判日承担更重的账目，并使你的惩罚不可忍受。因为报复是如此的邪恶，以至于它实际上撤销了天主的仁慈，并取消了已经赐予的无数罪过的赦免。因为那领受了万塔冷通债务豁免的人，在仅凭请求就获得了如此大的恩惠之后，却向自己的同伴索要一百个德纳，也就是说，要求同伴为自己的过犯赔补，他对同伴的严厉招致了自己的定罪；正是为此，他毫无例外地被交给刑役，遭受折磨，并被迫偿还那万塔冷通；他不被允许有任何借口或辩护，而是承受了最严厉的惩罚，被命令交出那曾因天主的慈爱而豁免的全部债务。\BibleRef{玛 18:23}那么，请问，财富之所以被你们如此热切地追求，难道就是因为它如此轻易地诱使人犯这种罪吗？不，绝非如此；这正是你们应当憎恶它如同一个充满无数杀戮的仇敌和对手的原因。但有人会说，贫穷使人不满，也常常使人说出亵渎的话，并屈就于卑劣的行为。不是贫穷导致这些，而是心灵的狭隘：因为拉匝禄也是穷人，是的，非常穷；除了贫穷，他还遭受病弱，这比任何形式的贫穷更为痛苦，并且使贫穷更加难以忍受；除了病弱，还有完全缺乏保护者，以及难以找到供应他所需的人，这加重了贫穷和病弱的痛苦。因为每一样痛苦本身都是痛苦的，但当没有人去照顾受苦者的需要时，痛苦就变得更大，火焰更炽烈，苦难更深重，风暴更猛烈，波浪更强，火炉更热。若仔细考察，除了这些还有第四重考验——近旁那富翁的冷漠与奢华。若你找到第五样作为火焰燃料的事物，你会清楚地看到他也被它所困。因为不仅那富翁生活奢侈，而且一天两次、三次，甚至多次他看见那穷人：因为穷人被放在他的门口，成为悲惨苦难的可怖景象，单是看见就足以软化铁石心肠：然而即便如此，也未使那无情的人帮助这穷困；但他摆开奢华的筵席，花环环绕的酒杯，丰盛的纯酒，大群厨师、食客和谄媚者从清晨开始，还有成队的歌手、司酒者和弄臣；他终日沉溺于各种放纵、酗酒和暴食，以及服饰、宴乐等许多事物。虽然他每天都看见那穷人因饥饿和极重的病弱，以及众多疮伤的压迫、贫困和由此产生的苦难而痛苦，却从未想到他：但食客和谄媚者却得到超过所需的滋养；而穷人，如此贫穷，被如此多的痛苦所包围，甚至得不到那桌上掉落的碎屑，尽管他极其渴望；然而这一切都没有伤害他，他没有说出一句苦毒的话，没有发出亵渎的言语；而是像一块在极热中炼净而更加明亮的金子，同样，他虽被这些苦难压迫，却超越了一切，也超越了这些苦难常常带来的躁动。因为一般来说，穷人看见富人时，往往被嫉妒消耗，被恶意恨恶折磨，认为生命不值得活，即使他们饮食充足，有人照顾；那么，这个穷人若不是非常明智和高贵，他的处境会怎样呢？他比所有其他穷人都穷，不仅贫穷，而且病弱，无人保护或安慰，躺在城中如同在遥远的荒漠，被饥饿折磨，看见一切好东西如同从泉源涌出般倾倒给富翁，却得不到任何人间的安慰，反而躺卧作为狗舌的常餐，因为他身体虚弱破碎，无法驱赶它们。你明白吗？那不自害的人不受任何损害。因为我要再次提出同样的论点。\newline{}\par

11. 因为这英雄的肉体病弱对他有何伤害？或是缺乏保护者？或是狗的到来？或是那富翁的恶邻？或是后者的极度奢侈、傲慢与骄横？这些是否削弱了他为美德争战的力量？是否毁了他的坚毅？他丝毫未被伤害，反而那众多的苦难和富翁的残酷增强了他的力量，成为他无限胜利冠冕的保证，增加他赏报的手段，扩充他报酬的应许。因为他得冠冕不只是因为他的贫穷、饥饿、疮伤或狗舔舐它们，而是因为：他有那样一个富翁为邻，天天被他看见，却永远被忽视，他勇敢而极其坚毅地忍受了这考验——这考验为贫穷、病弱和孤独之火添加了不小的火焰，实际上是极大的火焰。\newline{}\par

请问，有福的\Person{保禄}的情况如何？因为我没有什么可以阻止我再次提及他。难道他没有经历无数的考验风暴吗？他因此受到了什么伤害？他岂不更是因此获得了胜利的冠冕——因为他遭受饥饿，因为备受寒冷和赤身之苦，因为时常受鞭笞，因为被石头砸，因为被抛入海中？但有人说：\Quote{他是\Person{保禄}，并蒙\Person{基督}召唤。}然而\Person{犹达斯}也是十二人之一，他也蒙\Person{基督}召唤；但无论是十二人的身份还是他的召唤都没有使他获益，因为他没有一颗倾向美德的心。而\Person{保禄}虽然与饥饿斗争，无力获得必需品，每日承受如此巨大的苦难，却以极大的热诚追求通往天堂的道路；而\Person{犹达斯}虽然在他之前蒙召，享有与他相同的优势，领受了基督徒生活的最高形式，参与了圣桌和那最可敬畏的神圣筵席，并领受了能复活死人、洁净癞病人、驱逐魔鬼的恩宠，常常听讲关于贫穷的言论，并与\Person{基督}本人共度如此长的时间，还被委托管理穷人的钱财，好借此平息他的贪欲（因为他是个贼），即便如此也没有变得更好，尽管他蒙受了如此大的屈尊俯就。因为\Person{基督}知道他贪财，并注定因贪爱钱财而灭亡，所以当时不仅不为此惩罚他，反而为了缓和其贪欲，将穷人的钱财委托给他，使他能有某种止息贪婪的手段，以免落入那可怕的罪恶深渊，以此用较小的恶预先阻止了更大的恶。\par

12. 因此，在任何情况下，没有人能伤害一个不愿伤害自己的人；但如果一个人不愿节制，不愿借助自身资源帮助自己，那么没有人能使他获益。因此，圣经中那奇妙的历史，如同在一幅高大、宽阔、广袤的图画中，描绘了古人的生平，将叙述从\Person{亚当}延伸到\Person{基督}的来临；并向你展示了那些跌倒的和在竞赛中得胜冠冕的人，为要通过所有榜样教导你：即使整个世界向一个人燃起猛烈的战争，也没有人能伤害那不被自己伤害的人。因为不是环境的压力、季节的变化、权贵者的侮辱、如暴风雪般包围你的阴谋、一大群灾祸、或人类所遭受的种种不幸的混杂集合，能稍微动摇那勇敢、节制、警醒的人；相反，那懒惰、怠惰、出卖自己的人，即使有无数帮助，也不能变得更好。这一点至少由两个人的比喻向我们显明了：一人把房子建在磐石上，另一人建在沙土上：这并不是要我们想到沙土和磐石，或是石造建筑和屋顶，或是河流、雨水、狂风击打房屋，而是要从这些事中提取美德和恶习作为意义，并从中领悟：没有人能伤害一个不伤害自己的人。因此，无论是猛烈倾泻的雨水，还是猛烈冲撞的急流，还是猛烈冲击的狂风，都丝毫未能摇动那第一座房屋；它却安然无恙，纹丝不动：好让你明白，任何考验都不能动摇那不出卖自己的人。但另一个人的房屋轻易被冲毁，不是因为考验的猛烈（因为如果这样，第一座房屋也会遭遇同样命运），而是因为他自己的愚昧；它倒塌不是因为风吹，而是因为建在沙土上，即建在懈怠和不义之上。因为在暴风雨冲击之前，它已经软弱、摇摇欲坠。因为这类建筑，即使无人施加压力，也会自行碎裂，地基处处下陷、松动。正如蜘蛛网即便不被拉扯也会自然裂开，但金刚石即使被击打也岿然不动：同样，那些不伤害自己的人，即使受到无数打击，反而变得更加坚强；但那些出卖自己的人，即使无人骚扰，也会自行跌倒、崩溃、灭亡。因为\Person{犹达斯}就是这样灭亡的：他不仅没有受到此类考验的攻击，反而实际上享受了许多帮助的益处。\par

13. 你愿意我用整个民族的例子来说明这个论点吗？对犹太民族曾赐予了何等伟大的眷顾！难道不是整个有形受造界都为了服侍他们而安排的吗？难道一种全新而奇异的生活方式不是在他们中间被引入的？因为他们不必去市场采购，无需付任何代价便享有那些用钱才能买到的东西；他们不犁沟，不拉犁，不耙地，不撒种，无需雨水和风、四季更替、阳光、月相、气候以及诸如此类的一切；他们不预备打谷场，不打谷，不用簸箕分离谷糠，不推磨，不建炉灶，不把柴火带进屋里，不需要面包师傅的手艺，不用铲子，不磨镰刀，不需要其他技艺——我指的是纺织、建筑或制鞋；但天主的圣言就是他们的一切。他们随时备好筵席，无需任何辛劳。因为玛纳的性质正是如此：它总是新鲜、清新，不给他们带来任何麻烦，也不使他们因劳动而疲惫。他们的衣服、鞋子，甚至他们的身体，都忘记了自然的软弱：前者在如此漫长的岁月里没有磨损，他们的脚也没有肿胀，尽管他们长途跋涉。医生、医药以及所有关于那类技艺的操心，在他们中间完全无人提及；各种软弱就这样被彻底驱逐了：因为经上记载：\BibleQuote{他领他们出来，携带着金银，各支派中没有人残跛。} 又如那些已离开这个世界、被移植到另一个更美好世界的人，他们这样吃喝，炎热的阳光也不灼伤他们的头，因为云彩为他们挡住了炙热的光线，环绕在他们四周，像一座便携的庇护所，保护着全体人民。夜间他们也不需要火炬驱散黑暗，因为他们有火柱，一种说不出的光，满足两种需要：一是照亮，二是引导他们行路的方向；因为火柱不仅发光，而且比任何人向导更确定地引导着那无数的队伍在旷野中前行。他们不仅旱地行，还像在干地上一样在海中走；他们大胆地试验自然律，踏在那愤怒的海上，如踏在坚硬抵抗的岩石上一样行进；的确，当他们把脚放上去时，那元素变得像坚实的土地，以及平缓倾斜的平原和田野；但当它接纳他们的敌人时，它又按海的本性行事：对以色列人来说，它成了战车，对他们的敌人却成了坟墓；轻松地运载前者过去，却猛烈地淹死后者的。那无序的水流展现出理性而高度明智之人的良好秩序与顺从，一时扮演守护者的角色，一时扮演行刑者的角色，并在同一天展现出这些对立面。关于那涌出溪流的岩石，又该说什么呢？关于那布满整个地面的鸟尸之云呢？关于\Place{埃及}的奇迹呢？关于旷野的神迹呢？关于凯旋与不流血的胜利呢？因为他们降服了反对者，仿佛是过节而非作战；他们不用武器就战胜了自己的主人；离开\Place{埃及}后，他们用歌唱和音乐克服了与他们争战的人；他们所行的更像一个节庆而非战役，更像一次宗教仪式而非战斗。因为这一切奇事不仅仅是为了满足他们的需要，也是为了让他们更准确地保存\Person{梅瑟}所传授的关于认识天主的教义；四面八方都传出宣告他们主人临在的声音。大海高声宣告此事，因为它成为他们行路的大道，然后又变回海；\Place{尼罗河}的水也发出这声音，因为它们变成了血的性质；青蛙、大群的蝗虫、毛虫与霉也向全体人民宣告同样的事；旷野中的奇迹——玛纳、火柱、云彩、鹌鹑以及所有其他事件——都成为他们永不磨灭的书卷和文字，每天在他们的记忆中回响，在他们的心中回荡。然而，在如此伟大而卓越的眷顾之后，在这一切说不出的恩惠之后，在如此强大的奇迹之后，在无以言表的眷顾之后，在持续不断的教导之后，在藉言语训诫和藉行为劝戒之后，在光荣的胜利之后，在非凡的凯旋之后，在丰盛的食物供应之后，在充足的活水生产之后，在人类眼前所赋予他们的不可言喻的荣耀之后，他们却忘恩负义、愚昧无知地崇拜了一只牛犊，向一个牛头致敬，即使天主在\Place{埃及}的恩惠的纪念在他们心中还是新鲜的，他们仍然实际享有更多的恩惠。\par

14. 但是尼尼微人，尽管是一个野蛮的外邦民族，从未领受过这些恩惠，无论大小，既无言语，也无奇事，亦无作为，当他们看见一个从海难中获救的人——这人先前从未与他们交往，而是初次出现——进入他们的城，说：\BibleQuote{再过三天，尼尼微就要倾覆了。}\BibleRef{纳 3:4} 仅仅因这句话的声响就如此悔改和更新，放下从前的邪恶，沿着悔改之路向德行迈进，致使天主的判决被撤销，阻止了那威胁他们城市的动乱，避开了来自上天的愤怒，并从一切邪恶中得了救。\Quote{因为，}我们读到，\BibleQuote{天主看见每个人都离开了自己的邪道，归向了主。}\BibleRef{纳 3:10} 怎样离开？我问。虽然他们的邪恶甚大，他们的不义不可言说，他们的道德创伤难以治愈，这明显地由先知指出来，当他曾说：\BibleQuote{他们的邪恶甚至上达于天。}\BibleRef{纳 1:2} 以地域的距离表明他们邪恶的严重性；然而如此巨大的不义，堆积到如此高度，以致上达于天，这一切在三天之内，通过片刻时间，借着从一个人——一个陌生、遭遇海难的人——口中所说的几句话，他们如此彻底地废除、消除、除去，以至于幸而听到声明：\BibleQuote{天主看见每个人都离开了自己的邪道，天主就后悔了，不再降灾。}\BibleRef{纳 3:10} 你看见了吗？那节制和警醒的人不仅不被人所伤，反而能扭转上天的愤怒。而那背叛自己、自害己身的人，即使领受无数恩惠，也得不到多大益处。至少，犹太人并未因那些伟大的奇迹而受益，同样，尼尼微人也并未因没有分享那些奇迹而受害；但因他们内心向善，抓住了一个小小的机会就变得更好，尽管是野蛮的外邦人，对一切神圣启示一无所知，且住在远离巴勒斯坦的地方。\par

15. 我又问：那\Quote{三个少年}的德行是否被包围他们的患难所败坏？他们尚年幼，是少年，未及成熟年龄，是否承受了被掳的沉重痛苦？他们不是要长途跋涉离开家园，抵达异国后，与祖国、家园、圣殿、祭坛、祭献、供物、奠祭，甚至咏唱圣咏绝缘吗？因为不仅他们被剥夺了家园，而且因此也失去了许多敬拜的形式。他们不是被交在野蛮人——更像是豺狼而非人——的手中吗？而最痛苦的不幸是，当他们被放逐到如此遥远野蛮的国度，忍受如此沉重的俘虏生活时，他们没有导师，没有先知，没有领袖？因为经上记载：\Quote{因为没有领袖，没有先知，没有统治者，也没有在祢面前献祭并找到怜悯的地方。} 不仅如此，他们被投入王宫，如同被抛在悬崖峭壁上，面对一片岩石和暗礁的海洋，被迫在没有领航员、没有信号员、没有水手、没有帆的情况下航行那愤怒的海洋；他们被关在王庭，如同在监狱里。因为他们认识属神的智慧，超脱世俗的事物，藐视一切人类的傲慢，使灵魂的翅膀向上飞扬，所以他们把在那里寄居视为加重的痛苦。因为若他们不在宫中，而住在私人家里，他们会享有更大的自由；但被投进那个监狱（因为他们认为王宫的辉煌不比监狱好，不比岩石和悬崖之地更安全），他们即刻陷入残酷的窘境。因为国王命令他们共享自己的宴席，一个奢侈、不洁和亵渎的宴席，这是禁止他们的事，似乎比死亡更可怕；而他们孤独一人，像羔羊被狼群包围。他们被迫在要么被饥饿吞噬或更确切地被处决，要么品尝禁食之间做出选择。那么，这些少年——孤独、被掳、陌生、是那些命令者的奴隶——做了什么？他们不认为这种困境或拥有国家者绝对权力足以证明他们的顺从；他们用尽一切手段和方法来避免犯罪，尽管他们在各方面都被遗弃。因为他们无法用金钱影响人：作为俘虏，他们怎能？无法凭借友谊和社交：作为陌生人，他们怎能？无法凭借任何权力的行使来胜过他们：作为奴隶，怎么可能？无法凭借人数优势制服他们：只有三人，怎能？因此，他们去找那拥有必要权力的太监，用他们的论据说服了他。因为他们看见他恐惧战栗，为自身安全极度焦虑，死亡恐惧撕裂他的灵魂，难以忍受：\Quote{因为我害怕我的主上国王，}他说，\Quote{恐怕他看见你们的面容比你们同岁的少年憔悴，你们就会使我的头在国王面前不保。}\BibleRef{达 1:10} 他们使他从恐惧中解脱后，说服他赐予他们恩惠。因为他们将全部力量投入工作，天主也从那时起贡献了他的力量。因为他们在那些将获得赏报的事上成就的，并非只是天主的作为，开始和起点来自于他们自己的意志，当他们显明了那意志是高贵和勇敢时，他们为自己赢得了天主的帮助，从而实现了目标。\par

16. 你难道看不出，若一个人不伤害自己，就没有人能伤害他吗？请看：青春、被掳、贫困、流落异乡、孤独、缺乏保护、严峻的命令、面对死亡的巨大恐惧侵袭太监的心、贫穷、人数稀少、住在蛮族中间、以仇敌为主人、被交到君王手中、与所有亲人分离、远离司铎和先知、以及所有照顾他们的人、奠祭和祭献的停止、圣殿和圣咏的丧失——所有这些都未能伤害他们；相反，那时他们所得的荣耀更胜于在家乡享受这一切的时候。在他们首先完成这项任务，用光荣的胜利冠冕装饰了额头，即使在异乡也遵守了律法，践踏了暴君的命令，克服了对复仇者的恐惧，却没有受到任何伤害，仿佛他们安静地在家乡享受我所提的一切恩惠一样；在他们这样无畏地完成工作后，他们又被召唤参加其他的竞赛。他们仍是同样的人，遭受了比先前更严峻的考验：火窑被点燃，他们与君王一起面对蛮族大军；整个\Place{波斯}军队都行动起来，一切能施加欺骗或强迫的手段都被设计出来：各种音乐、各种刑罚、威胁，以及他们四面所见令人惊恐的景象，他们所闻的话语比所见更令人恐惧；然而，由于他们没有背叛自己，而是发挥自己的力量，他们从未遭受任何损害；反而为自己赢得了比先前更光荣的胜利冠冕。因为\Person{拿步高}捆绑他们，将他们投入火窑，却没有烧毁他们，反而使他们受益，使他们更加显赫。虽然他们失去了圣殿（我要重复先前的话）和祭台、祖国、司铎、先知，虽然他们身在异邦蛮族之地，在火窑的中央，被那强大的军队包围，行这事的君王本人也在观看，他们却立下了光荣的凯旋，赢得了显著的胜利，唱出了那首从那时直到今日在普世传唱、且将传给后世的可赞可颂的非凡颂歌。\par

因此，当一个人不伤害自己时，他就不可能被他人伤害：因为我要不断地重复这句话。因为被掳、捆绑、孤独、失去祖国和所有亲人、死亡、焚烧、大军、残暴的君王——这一切都不能损害那三个青年（他们是俘虏、奴仆、身在异乡的陌生人）的内在德行，反而使敌人的攻击成为他们更大信心的机会。那么，有什么能伤害节制的人呢？没有，即使整个世界上的人都武装起来反对他。但也许有人说：在他们的情况下，天主站在他们身边，把他们从火焰中救了出来。确实，祂这么做了；如果你尽你所能履行你的本分，天主所赐的帮助必定随之而来。\par

17. 然而，我之所以钦佩这些青年，称他们为有福的、可羡慕的，并不是因为他们践踏火焰，战胜了火的力量；而是因为他们为了真道而被捆绑、投入火窑，被交于烈火。因为正是这一点构成了他们胜利的圆满，胜利的花环在他们被投入火窑之时就已放在他们的额上，并且在事件的结果之前，从他们被带到王面前时所说的那些充满勇气和直言的话开始，就已经为他们编织了。他们说：\BibleQuote{我们不需要在这件事上回答你；因为我们所侍奉的天主在天上，祂能拯救我们脱离这燃烧的火窑；祂必从你手中拯救我们，大王。即便不然，大王，你也当知道，我们决不侍奉你的神，也不敬拜你所立的金像。}\BibleRef{达 3:16}说出这些话之后，我就宣布他们为得胜者；说了这些话之后，他们抓住了胜利的奖赏，冲向殉道的光荣桂冠，用言语所作的宣认继而以行动所作的宣认跟上。但如果当他们被投入火中时，火尊重他们的身体，解开了他们的束缚，允许他们毫不惧怕地下去，并且忘记了它的自然力量，以至于火窑变成了清凉的水泉，这奇迹是天主恩宠和神圣奇能的效果。然而，这些英雄甚至在这一切发生之前，一旦踏入火焰，就已经竖立了他们的胜利纪念碑，赢得了胜利，戴上了冠冕，并在天上和地上被宣布为得胜者，就他们而言，他们的荣誉没有任何缺失。那么，你对此有何话说？你被流放，被逐出祖国？看，他们也如此。你遭受了被掳，成为了蛮族主人的仆人？同样，你会发现这些人也遭遇了这些。但你身边没有人来管理你的状况，也没有人指导你或劝诫你？这些人都缺乏这种关怀。或者你被捆绑、被焚烧、被处死？因为你不能告诉我比这些更痛苦的事了。然而，看！这些人都经历了这一切，且因每一件事而更加光荣，更加显赫，并增加了他们在天上的宝藏。那些犹太人虽然拥有圣殿、祭坛、约柜、革鲁宾、赎罪盖、帐幔，无数的司祭，每日的礼仪，早晚的祭献，不断听到先知的声音（无论是活着的还是已逝的）在他们耳边响起，随身携带着对在埃及和旷野所行奇事的回忆，以及所有其余的事，并将这些事情的故事在手中翻阅，写在门框上，并在那时享受了许多超自然的能力和其他各种帮助，却毫无益处，反而受损，甚至在圣殿本身设立了偶像，在树下献祭了他们的儿女，在巴勒斯坦几乎每一个地方都献上了那些不合法且可咒骂的祭品，并犯下了无数更可怕的行为。但这些人在野蛮敌对之地，身处暴君之家，被剥夺了我所说的所有关怀，被带去处死，遭受焚烧，不仅没有受到丝毫伤害，反而变得更加显赫。知道了这些，并从圣神所默感的圣经中收集类似的例子（因为可以在许多其他人物身上找到许多这样的例子），我们就认为无论是时势或事件的困难，还是强迫和暴力，或是权贵的专横权力，都不能为我们犯罪提供充分的借口。我现在要重复我在开头所说的话来结束我的讲论：如果有人受到伤害和损害，他肯定是在自己手里遭受这一切，而不是在别人手里，即使有无数的众人伤害和侮辱他；所以，如果他没有在自己手里遭受这一切，那么所有居住在普天之下和海中的受造物，即使联合起来攻击他，也不能伤害一个在主内警醒和清醒的人。因此，我恳求你们，要时常清醒和警醒，勇敢地忍受一切痛苦，好使我们能在我们的主耶稣基督内获得那些永恒而纯洁的福乐，愿光荣和权能归于祂，从今世直到永远。阿们。\par

\SourceRecord{Translated by W.R.W. Stephens.；Nicene and Post-Nicene Fathers, First Series；Vol. 9.；Edited by Philip Schaff.；Buffalo, NY: Christian Literature Publishing Co.,；1889.；<http://www.newadvent.org/fathers/1902.htm>.}
